\begin{proposition}[the canonical ideal need not be idempotent]
\label{prop:canonical-nonidempotent}
There is a finite distributive, commutative, entire idempotent semiring
$Q$ and a finite free $Q$-semimodule $L$ for which the canonical ideal
of $G_Q(\Z)$ satisfies $\mathcal J_L^2\ne\mathcal J_L$.
\end{proposition}
\begin{proof}
Let $Q=\{0,a,b,c,1\}$, ordered by
$0<a<b<1$ and $0<a<c<1$, with $b,c$ incomparable.
This is the ideal lattice of a three-point poset with one least point
and two incomparable larger points. Addition is join and multiplication
is meet. Thus $b\vee c=1$, $b\wedge c=a$, and every two nonzero
members have nonzero meet. In $T=G_Q(\Z)$ only the top coefficient
group is nonzero, and it is $\Z$.

Take $L=Q^2$, $o=(0,0)$ and $k=(b,c)$. Here $\mathcal A_L$ is exactly
the $\Z$-linear incidence category of $L$: every proper-support scalar
arrow is zero, the top scalar $n$ is $n1_l$, and the remaining relations
are those of the transition arrows. Conversely these assignments in the
incidence category satisfy every defining relation, so no further
relations are imposed.

Every $Q$-linear functional has the form
\[
 \phi(u,v)=(u\wedge d_1)\vee(v\wedge d_2).
\]
A nonzero canonical generator at target $k$ requires $\phi(k)=1$.
Both $d_1,d_2$ must then be nonzero, because its individual summands
are bounded by $b<1$ and $c<1$. If $\phi$ is admissible at
$l=(l_1,l_2)$, evaluation at $(1,0)$ and $(0,1)$ gives
$d_1\wedge l_2=d_2\wedge l_1=0$. Entirety forces $l=o$.
Consequently $\mathcal J_L(l,k)=0$ for $l\ne o$.
At $o$, the functional $\phi(u,v)=u\vee v$ is admissible and has
$\phi(k)=1$, so its coefficient $1$ gives
$\mathcal J_L(o,k)=\Z t_{o,k}$. On the other hand
$\mathcal J_L(o,o)=0$, because every functional preserves zero and
$F(0)=0$. In every composite contributing to
$\mathcal J_L^2(o,k)$, the intermediate object must therefore be $o$,
and the remaining factor is zero. Hence
\[
 \mathcal J_L(o,k)=\Z,\qquad \mathcal J_L^2(o,k)=0.
\]
This proves the assertion. It does not contradict the criterion:
$\mathcal J_LG=G$ implies $\mathcal J_L^nG=G$ for every $n\geq1$,
without requiring equality of the ideals themselves.
\end{proof}

\begin{proposition}[every ordinary ideal occurs]
\label{prop:canonical-arbitrary-ideal}
Let $R$ be any unital ring and $I$ any two-sided ideal, without a
unitality assumption on $I$. There is a two-idempotent inverse semiring
$T$ and a finite projective support $L$ for which
$\mathcal A_L=R$ and $\mathcal J_L=I$.
A semimodule on this support is finite projective over $T$ exactly when
its underlying $R$-module $P$ is finite projective and $IP=P$.
In particular, powers of the canonical ideal need not stabilize.
\end{proposition}
\begin{proof}
Take the disjoint union of a lower copy of $I$ and an upper copy of $R$.
Add within each group in the usual way; a mixed sum lies in the upper
copy and uses $I\hookrightarrow R$. A product lies in the lower copy
if either factor does, and otherwise in the upper copy, with its
coefficient given by ring multiplication. The two-sided ideal property
makes this well defined. Associativity and distributivity hold on
coefficients by the ring identities and on tags by the meet and join
laws of $\B$. The lower zero is the absorbing additive identity, the
upper $1_R$ is the multiplicative identity, and additive inverse is
negation within each group. Thus this is the injective two-idempotent
construction with $F(0)=I$, $F(1)=R$ and transition the inclusion.

Let $L=\{o\}$ be the zero $\B$-semimodule, which is finite projective.
Its action category has one object. The upper scalar arrows give $R$;
the coefficient transition relation identifies each lower scalar
$i\in I$ with that same element of $R$. Conversely the ordinary
$R$-action satisfies all defining relations. Hence there are no further
relations and $\mathcal A_L$ is the one-object category $R$.
The only $\B$-linear functional $L\to\B$ is zero. It is admissible,
and its canonical generators are precisely $F(0)=I$. Therefore
$\mathcal J_L=I$. The projectivity assertion is now
Theorem~\ref{thm:canonical-ideal}. Taking $R=\Z$ and $I=2\Z$ gives
$\mathcal J_L^n=2^n\Z$, which is strictly decreasing for every $n$.
\end{proof}
These examples separate the canonical support ideal from the trace
ideal of a particular finite projective. They also show that arbitrary
ordinary ideal-generated projective modules occur in the residual
classification problem.
